\documentclass[10pt,a4paper]{amsart}
\usepackage[margin=19mm]{geometry}
\usepackage{amsmath,amssymb,mathtools}
\usepackage[scr=rsfs,cal=euler]{mathalfa}
\usepackage{xfrac}
\usepackage[unicode,hidelinks,bookmarks=false]{hyperref}
\newcommand{\ie}{i.e.\ }
\newcommand{\cf}{cf.\ }
\newcommand{\resp}{respectively\ }
\newcommand{\Subsection}{\subsection}
\providecommand{\nobreakdash}{\nobreak\hbox{-}}
\setlength{\emergencystretch}{2em}
\multlinegap0pt
\usepackage[shortlabels]{enumitem}
\usepackage{amssymb}
\usepackage{mathtools}
\newtheorem{corollary}{Corollary}
\newtheorem{theorem}{Theorem}
\newtheorem{lemma}{Lemma}
\newcommand{\abs}[1]{| #1 |}
\DeclareMathOperator{\dist}{dist}
\newcommand{\inner}[2]{\langle #1, #2 \rangle}
\newcommand{\norm}[1]{\| #1 \|}
\title{Anderson acceleration of the proximal point method: the exact adaptive minimax, a spectral phase transition, and optimal safeguarding}
\author{Zheng Jia, Yekini Shehu, Yonghong Yao}
\date{Source: arXiv:2607.24643v1 (CC BY 4.0)}
\begin{document}
\maketitle

\noindent\emph{Source and licence.} Anderson acceleration of the proximal point method: the exact adaptive minimax, a spectral phase transition, and optimal safeguarding. Version arXiv:2607.24643v1 (CC BY 4.0), distributed by arXiv under the Creative Commons Attribution 4.0 International licence, http://creativecommons.org/licenses/by/4.0/, as recorded in the arXiv licence metadata for this exact version and shown on the abstract page https://arxiv.org/abs/2607.24643v1. Full licence notice: \texttt{licenses/arxiv-2607.24643-CC-BY-4.0.txt}.\par\smallskip

\noindent\textit{Editorial scope.} Counted unit: the article's theorem thm:growth (source0.tex lines 1333--1350). The article's complete original proof of it is delivered with it (source0.tex lines 1351--1378, one \texttt{proof} environment of 28 lines). No mathematics is written, repaired or re-derived: the statement and the proof are the article's own lines, in the article's own notation, and no retained line is substituted or re-flowed.  Retained uncounted as the source-local context the proof needs: lines 430--434; lines 983--995; lines 1185--1200; lines 1324--1327; lines 1831--1831; lines 1833--1840.  Nothing was omitted from any retained span, so the package carries no omission record.  No retained line contains a citation, so the wrapper carries no bibliography and no bibliography entry.  No retained line contains a figure environment, an image-inclusion command or drawing code, so no image is delivered and the package declares no asset dependency.  Editorial formatting only, in addition: an AMS \texttt{amsart} preamble that restates the article's own declarations and macros used by the retained lines, namely the environments \texttt{corollary}, \texttt{theorem}, \texttt{lemma} on separate counters, together with the standard packages those lines need.  The retained environments print in this document as Corollary 1, Theorem 1, Lemma 1 in order of appearance, whereas the article prints the same items with the numbering of its own full document.  The complete native source of the article as distributed in the arXiv e-print for this exact version is bundled unchanged as \texttt{sources/206129/source0.tex}.
\begin{corollary}[Linear safety: no safeguard needed]\label{cor:linsafe}
On every linear maximal monotone problem, the AA-PPM residual sequence is
nonincreasing, $\norm{r_{k+1}}\le\norm{r_k}$ for all $k$, with any memory
$m$. Safeguarding is therefore a purely \emph{nonlinear} phenomenon.
\end{corollary}

\begin{theorem}[Escape: Jackson kernel on floored spectra]\label{thm:jackson}
Let $N=\lfloor(K+2)/2\rfloor$ and
$q_J(u)=\big(\frac1N\sum_{j=0}^{N-1}u^j\big)^2$ (the squared
Fej\'er/Jackson kernel, $\deg q_J=2N-2\le K$, $q_J(1)=1$). For every
rotation spectrum with floor $s=\dist(1,\sigma(L))$ and every $y_0$,
\begin{equation}\label{eq:jackson}
\norm{r(y_K)}\ \le\ \frac{4\,d_0}{N^2 s}
\ \le\ \frac{16\,d_0}{(K+1)^2 s}.
\end{equation}
In particular $\norm{r(y_K)}=o(d_0/K)$ whenever $sK\to\infty$, and
$q_J$ improves on the universal envelope $1/(K+1)$ as soon as
$s\gtrsim 16/(K+1)$.
\end{theorem}

\begin{theorem}[Certificate and rate inheritance]\label{thm:cert}
Algorithm~1 uses exactly $2K+1$ resolvent evaluations in $K$ iterations, and
\begin{enumerate}[label=(\roman*)]
\item \emph{PPM floor}: $\norm{r(b_K)}\le\norm{r^c_K}\le
      d_0/\sqrt{K+1}$ for all $K$; i.e.\ the certified output is never
      worse than PPM run at half the oracle budget, and inherits every
      rate PPM has under additional structure (strong monotonicity,
      growth conditions, etc.);
\item \emph{Acceleration side}: along every acceptance run
      $[k_0,k]$, $\norm{r(b_k)}\le(1-\delta)^{k-k_0}\norm{r^c_{k_0}}$;
      on problems where raw AA converges (super)linearly---e.g.\ the
      floored-spectrum regime of Theorem~\ref{thm:jackson}---the output
      inherits the accelerated rate with the floor of (i) intact;
\item \emph{Monotonicity}: $\norm{r(b_{k+1})}\le\norm{r(b_k)}$.
\end{enumerate}
\end{theorem}

\begin{equation}\label{eq:growth}
\inner{w}{x-\Pi_Zx}\ \ge\ \mu\, d(x)^q
\qquad\forall\, w\in Mx .
\end{equation}

\section{A sequence lemma}\label{app:seq}

\begin{lemma}\label{lem:seq}
Let $u_k\ge0$ satisfy $u_k-u_{k+1}\ge c\,u_{k+1}^{\alpha}$ with $c>0$,
$\alpha>1$. Then
\[
u_k\ \le\ \big(c(\alpha-1)\,k\big)^{-1/(\alpha-1)}
\qquad\text{for all }k\ge1 .
\]
\end{lemma}

\begin{theorem}\label{thm:growth}
Let $M$ be maximal monotone satisfying \eqref{eq:growth}, and let
$y_{k+1}=Jy_k$ (PPM), $d_k=d(y_k)$.
\begin{enumerate}[label=(\alph*)]
\item $q>2$: $d_k=O(k^{-1/(q-2)})$; if moreover
      $\dist(0,Mx)\le L\,d(x)^{q-1}$ (automatic for $\partial f$ of
      H\"olderian $f$), then
      $\norm{r_k}=O(k^{-(q-1)/(q-2)})$. Both exponents are sharp:
      $f(x)=\abs{x}^q/q$ attains them.
\item $q=2$: linear convergence,
      $d_{k+1}^2\le d_k^2/(1+2\mu)$.
\item $q<2$: local superlinear convergence of order $1/(q-1)>1$:
      $d_{k+1}\le (2/\mu)^{1/(q-1)}\,d_k^{\,1/(q-1)}$ once
      $d_k$ is small enough.
\end{enumerate}
The same rates hold for Algorithm~1 (by Theorem~\ref{thm:cert}(i)) and for
raw AA-PPM on linear problems (by Corollary~\ref{cor:linsafe}).
\end{theorem}

\begin{proof}
The resolvent identity gives $w_{k+1}=y_k-y_{k+1}\in My_{k+1}$. Firm
nonexpansiveness of $J$ with $z=\Pi_Zy_{k+1}$ yields
\[
d_k^2-d_{k+1}^2
\ge \norm{w_{k+1}}^2+2\inner{w_{k+1}}{y_{k+1}-z}
\ge 2\mu\, d_{k+1}^{\,q},
\]
where the first inequality is
$\norm{y_k-z}^2=\norm{y_k-y_{k+1}}^2+\norm{y_{k+1}-z}^2
+2\inner{y_k-y_{k+1}}{y_{k+1}-z}$ and the second uses \eqref{eq:growth}.

(a) $u_k=d_k^2$ satisfies $u_k-u_{k+1}\ge 2\mu u_{k+1}^{q/2}$ with
$q/2>1$; Lemma~\ref{lem:seq} (Appendix~\ref{app:seq}) gives
$u_k=O(k^{-2/(q-2)})$. With $\dist(0,Mx)\le Ld(x)^{q-1}$,
$\norm{r_k}=\norm{w_{k+1}}\le L d_{k+1}^{\,q-1}=O(k^{-(q-1)/(q-2)})$.
Sharpness: for $f(x)=\abs{x}^q/q$ the resolvent solves
$x+\abs{x}^{q-2}x=y$, and a direct asymptotic expansion gives
$d_k=\Theta(k^{-1/(q-2)})$ and $\norm{r_k}=\Theta(k^{-(q-1)/(q-2)})$.

(b) $q=2$: $d_k^2-d_{k+1}^2\ge2\mu d_{k+1}^2$, i.e.\
$d_{k+1}^2\le d_k^2/(1+2\mu)$.

(c) $q<2$: from \eqref{eq:growth} and Cauchy--Schwarz,
$\norm{w_{k+1}}\ge\mu d_{k+1}^{\,q-1}$; but
$\norm{w_{k+1}}=\norm{y_k-y_{k+1}}\le d_k+d_{k+1}\le 2d_k$, so
$d_{k+1}\le(2/\mu)^{1/(q-1)}d_k^{1/(q-1)}$.
\end{proof}

\end{document}
